\chapter{Diskussion und Ausblick}
\label{chap:diskussion-und-ausblick}

\section{Diskussion der Ergebnisse}
\label{sec:diskussion-der-ergebnisse}

\subsection{Redundanz statt Quittierung}
\label{subsec:diskussion-redundanz}

Die zentrale Entwurfsentscheidung der Funkstrecke besteht darin, auf einen
Rückkanal zu verzichten und die fehlende Zustellgarantie stattdessen durch
mehrfache Aussendung desselben Datenstands zu kompensieren
(Abschnitt~\ref{subsec:drahtlose-uebertragung}). Diese Entscheidung war zum
Zeitpunkt des Entwurfs eine Annahme: die Messungen belegen nun ihre
Tragfähigkeit. Über eine Distanz von 100\,m mit Sichtverbindung bleibt der
Sequenzverlust bei 0,09\,\%, obwohl die Empfangsredundanz zu diesem Zeitpunkt
bereits von fünf auf 4,43 Kopien je Sequenzwert gesunken ist. Die Strecke
verliert also messbar an Qualität, ohne dass sich dies auf die Lichtausgabe
auswirkt, genau dies ist die Funktion, die der Redundanz zugedacht war.

Aus dieser Beobachtung folgt ein Gedanke, der über die Bewertung der einzelnen
Messreihe hinausreicht. Die Empfangsredundanz ist eine Zustandsgröße der
Funkstrecke, die ohne Rückkanal allein empfängerseitig bestimmbar ist. Die
Röhre kann zählen, wie viele Kopien eines Sequenzwerts sie erhalten hat, und
daraus auf die verbleibende Fehlerreserve schließen. Ein unidirektional
ausgelegtes System kann seinen eigenen Verbindungszustand somit kennen, ohne die
in Abschnitt~\ref{sec:abgrenzung} bewusst ausgeschlossene bidirektionale
Kommunikation einzuführen. Diese Eigenschaft ist die Voraussetzung für jede Form
der Selbstanpassung und wird in Abschnitt~\ref{subsec:ausblick-messungen} wieder
aufgegriffen.

Die Messungen erlauben darüber hinaus eine genauere Bestimmung dessen, wogegen
die Redundanz überhaupt schützt. Die fünf Kopien eines Sequenzwerts werden
innerhalb desselben 25-ms-Fensters ausgesendet
(Abschnitt~\ref{subsec:messaufbau}). Die Redundanz ist damit eine rein zeitliche
Wiederholung innerhalb eines sehr kurzen Intervalls. Sie wirkt folglich gegen
unkorrelierte Einzelverluste eine Kollision, eine kurzzeitige
Kanalbelegung, ein einzelnes gestörtes Paket, weil in diesen Fällen mindestens
eine der übrigen Kopien den Empfänger erreicht. Gegen korrelierte Störungen,
die länger als 25\,ms andauern, ist sie hingegen wirkungslos: Überschreitet die
Störung dieses Fenster, gehen alle fünf Kopien desselben Werts gemeinsam verloren.

Genau dieses Verhalten belegen beide Messreihen. Die Unterbrechung von 750\,ms im
Stoßzeitlauf umfasst 30 aufeinanderfolgende Sequenzwerte und damit 150 einzeln
ausgesendete Pakete. Ebenso wachsen mit zunehmender Distanz nicht nur die
Verlustraten, sondern vor allem die Längen der zusammenhängenden Verlustserien.
Der praktisch relevante Fehlerfall ist somit nicht der Einzelverlust, gegen den
das Verfahren ausgelegt ist, sondern der zeitlich ausgedehnte Ausfall. Daraus
folgt eine Einschränkung für jede Weiterentwicklung des Ansatzes: Eine bloße
Erhöhung der Kopienzahl innerhalb desselben Fensters verbessert die Robustheit
gegenüber solchen Ereignissen nicht.

Ebenso deutlich zeigen die Messungen die Grenze des Prinzips. Der Sequenzverlust
wächst nicht proportional zur Verschlechterung der Strecke, sondern bleibt lange
nahezu unverändert und springt anschließend sprunghaft an, sobald die Redundanz
gegen eins geht. Im verdeckten Lauf über 200\,m stehen im Mittel noch 1,61 Kopien
je Sequenzwert zur Verfügung, in einzelnen Fenstern genau eine und der Verlust
erreicht 36,29\,\%. Das Verfahren puffert also über einen weiten Bereich sehr
wirkungsvoll und versagt dann innerhalb einer kurzen Spanne. Für die Praxis
bedeutet dies, dass ein System, das nur den Sequenzverlust beobachtet, keine
Vorwarnzeit besitzt: Bis kurz vor dem Zusammenbruch der Verbindung erscheint die
Übertragung fehlerfrei.

\subsection{Broadcast ohne Rückkanal}
\label{subsec:diskussion-broadcast}

Der Verzicht auf Quittierung ist Teil einer umfassenderen Entscheidung für eine
verbindungslose Broadcast-Topologie. Ihr Gewinn liegt in der Einfachheit: Der
Sendeaufwand der Bridge ist unabhängig von der Anzahl der Empfänger, alle Röhren
erhalten denselben Datenstand innerhalb desselben Übertragungsfensters, und die
Adressierung erfolgt lokal am Gerät statt über ein Funkprotokoll
(Abschnitt~\ref{subsec:synchronisation-roehem}). Die Gleichzeitigkeit ist dabei
nicht nur eine Eigenschaft des Entwurfs, sondern messtechnisch bestätigt: In der
Zeitlupenaufnahme mit 120~Bildern je Sekunde wechseln acht gemeinsam betriebene
Röhren innerhalb desselben Einzelbildes
(Abschnitt~\ref{sec:gleichzeitigkeit}). Der Preis dafür ist ebenso klar:
Es existiert weder ein Zustellnachweis noch eine Möglichkeit zur Ferndiagnose. Ein
Ausfall einzelner Röhren wird nicht gemeldet, sondern erst optisch sichtbar.

Die Störfestigkeitsmessung erlaubt es, diesen Zielkonflikt genauer zu bewerten,
als es beim Entwurf möglich war. Das einzige praktisch relevante Ereignis des
Stoßzeitlaufs eine zusammenhängende Unterbrechung von 750\,ms lässt sich am
ehesten auf die Konkurrenz um den Übertragungskanal zurückführen und nicht auf
die fehlerhafte Zustellung einzelner Pakete: Sämtliche Kopien des betroffenen
Zeitabschnitts blieben aus, während Abstand und Aufbau unverändert waren
(Abschnitt~\ref{sec:stoerfestigkeit}). Ein unmittelbarer Nachweis ist mit dem
verwendeten Aufbau nicht möglich, da senderseitig nicht protokolliert wird, ob
ein Paket abgesetzt werden konnte. Unter dieser Deutung hätte ein
Quittierungsverfahren das Ereignis nicht verhindert. Es hätte den Ausfall
lediglich erkennbar gemacht und zugleich zusätzliche Sendevorgänge auf einem
bereits ausgelasteten Kanal erzeugt.
Die Entscheidung gegen den Rückkanal wird durch die Messung also nicht relativiert,
sondern gestützt allerdings mit der Einschränkung, dass die fehlende
Rückmeldung den Betreiber im Fehlerfall ohne Information lässt.

Der Zugriff auf das Medium erfolgt nach demselben CSMA/CA-Verfahren wie
bei WLAN, sodass die verfügbare Sendezeit mit der Fremdlast auf dem gewählten
Kanal konkurriert. Die Implementierung sieht hierfür drei gespreizte Kanäle vor
(Abschnitt~\ref{subsec:espnow}), überlässt die Auswahl jedoch dem Anwender. Der
Betrieb hängt damit an einer Konfigurationsentscheidung, die vor der Veranstaltung
zu treffen ist.

\subsection{Reichweite als Eigenschaft des Aufbaus}
\label{subsec:diskussion-abschattung}

Der aufschlussreichste Befund der Reichweitenmessung ist nicht die erreichte
Maximaldistanz, sondern der Vergleich beider Messreihen. Dieselbe Hardware, derselbe
Kanal und dieselbe Sendeleistung führen zu deutlich unterschiedlichen Ergebnissen,
je nachdem, ob die Sichtlinie frei ist oder durch ein massives Bauteil unterbrochen
wird. Die verdeckte Strecke verhält sich näherungsweise wie eine etwa doppelt so
lange freie Strecke (Abschnitt~\ref{subsec:reichweite-bewertung}). Reichweite ist
demnach keine Eigenschaft des Geräts, sondern eine Eigenschaft des Aufbaus. Für die
Anwendung folgt daraus eine unmittelbar umsetzbare Konsequenz: Die Positionierung
der Bridge auf einer möglichst freien Sichtachse und in erhöhter Lage wirkt
stärker auf die Verbindungsqualität als jede Steigerung der Sendeleistung, und eine
ungünstig platzierte Bridge lässt sich durch Reserven auf der Funkstrecke nicht
ausgleichen.

Dieser Befund verbindet die Arbeit mit der in
Abschnitt~\ref{subsec:forschung-kommunikation} beschriebenen Lösung von Yang et al.
Dort werden abgeschattete Empfänger über ein Kabel an einen erreichbaren Empfänger
angebunden, also durch das Wiedereinführen genau jener Verkabelung, deren
Vermeidung das erklärte Ziel des Ansatzes war. Die vorliegende Messreihe
quantifiziert erstmals, wie groß das zugrunde liegende Problem tatsächlich ist,
und zeigt zugleich, dass es im geforderten Einsatzbereich nicht zu diesem
Rückschritt zwingt: Bei der in Abschnitt~\ref{subsec:fa-kommunikation}
festgelegten Mindestreichweite von 50\,m bleibt die Übertragung auch bei
vollständig verdecktem Empfänger mit 2,08\,\% Sequenzverlust nutzbar.

Für den Fall größerer oder ungünstigerer Strecken bieten sich Gegenmaßnahmen an,
die den kabellosen Charakter des Systems erhalten. Naheliegend ist die Verteilung
mehrerer Bridges auf ein Veranstaltungsgelände, sodass jede Röhre eine kurze,
möglichst freie Strecke zu einem Sender besitzt. Die Broadcast-Topologie steht dem
nicht entgegen, erfordert jedoch die zuweise zu unterschiedlichen Kanälen 
inerhalb eines Sendebereichs um Interferenze zu vermeiden.
Ebenso wirksam ist die Ausrichtung und
Wahl der Antennen, die im vorliegenden Aufbau als externe Bauform ausgeführt sind
(Abschnitt~\ref{subsec:funk-antenne}) und damit ohne Eingriff in die Elektronik
getauscht werden können.

\subsection{Einordnung in den Stand der Technik}
\label{subsec:diskussion-einordnung}

Die in Abschnitt~\ref{sec:luecken-stand-der-technik} formulierte Lücke betraf
weniger die drahtlose Übertragung als solche als vielmehr den Grad der
Integration: Die betrachteten Arbeiten ersetzen jeweils nur die Übertragungsstrecke
zwischen Steuergerät und Leuchtmittel, während das Leuchtmittel selbst
kabelgebunden angesteuert bleibt. Das entwickelte System vereint Funkempfang,
Energieversorgung und Pixel Ansteuerung innerhalb der Röhre. Die
funktionale Verifikation in Abschnitt~\ref{sec:funktionale-tests} belegt, dass
mehrere unabhängig adressierte Röhren allein über die Funkstrecke betrieben
werden, und die Reichweitenmessung zeigt, dass dies über praxisrelevante
Distanzen trägt. Die Signalverkabelung entfällt damit vollständig und nicht nur
abschnittsweise.

Ein zweiter Beitrag betrifft die Bewertungsgrundlage. Die Leistungsbewertung in
\cite{yang2015real} beschränkt sich auf den Nachweis der Timing-Konformität.
Reichweite, Paketverlust und Verhalten unter Fremdfunklast bleiben unbetrachtet.
Die vorliegende Arbeit erhebt genau diese Größen und macht darüber hinaus die
Struktur der Verluste sichtbar. Dass ein Lauf mit 0,51\,\% mittlerem Verlust eine
zusammenhängende Unterbrechung von 750\,ms enthalten kann, während alle übrigen
Fenster nahezu fehlerfrei sind, ist für die Lichtwahrnehmung entscheidend und
bleibt bei einer Betrachtung von Mittelwerten unsichtbar. Die in
Abschnitt~\ref{subsec:metriken} eingeführte Trennung von Sequenzverlust und
Empfangsredundanz sowie die Auswertung der Verlustserien sind insofern
übertragbar auf andere funkgestützte Lichtsteuerungen und nicht auf das hier
entwickelte System beschränkt.

Gegenüber dem kommerziellen Premiumsegment bleiben Unterschiede bestehen, die
nicht auf die Übertragung, sondern auf den Funktionsumfang der Leuchte entfallen.
Die in Abschnitt~\ref{sec:kommerzielle-led-leuchtroehren} betrachtete
\textit{Titan Tube} bietet Akkubetrieb, eine Schutzart für den Außeneinsatz sowie
eine erweiterte Farbmischung mit zusätzlichen Mint- und Amber-Dioden, die eine
Anpassung des Weißabgleichs erlaubt. Das entwickelte System bleibt demgegenüber
netzgebunden und auf eine RGB-Mischung beschränkt. Die kabelgebundene
Pixelsteuerung ist auf der Bridge-Platine bereits vorbereitet, war jedoch
ausdrücklich nicht Gegenstand dieser Arbeit
(Abschnitt~\ref{subsec:bridge-pcb}).

Ein Vergleich der Kosten ist vor diesem Hintergrund nur mit Einschränkungen
aussagekräftig, verdeutlicht aber die Größenordnung. Den in
Abschnitt~\ref{sec:kommerzielle-led-leuchtroehren} genannten Endkundenpreisen von
etwa 800\,€ je Röhre im Premiumsegment und rund 400\,€ je Röhre bei
kostengünstigeren Systemen stehen Materialkosten von etwa 62\,€ je Röhre und rund
31\,€ je Bridge gegenüber (Abschnitt~\ref{subsec:gesamtkosten}). Diese Zahlen sind
jedoch nicht unmittelbar vergleichbar: Den die Bestückung
der Platinen, der 3D-Druck der Gehäuse, der mechanische Zusammenbau sowie die
Prüfung jeder Einheit, die im Eigenbau als Arbeitszeit anfallen und bei einem
Fertigprodukt im Preis enthalten sind fehlen. in diesen Preis. Ebenso wenig enthalten sind ein
EMV- und CE-Konformitätsnachweis und Gewährleistung.
Der Vergleich belegt daher nicht, dass das System um eine Größenordnung
günstiger ist, sondern dass die Materialkosten einer solchen Leuchte
deutlich unterhalb der Marktpreise liegen und der Eigenbau damit für Anwender
erschlossen wird, für die kommerzielle Systeme wirtschaftlich nicht in Frage
kommen. Eine weitergehende wirtschaftliche Betrachtung ist entsprechend der
Abgrenzung in Abschnitt~\ref{sec:abgrenzung} nicht Gegenstand dieser Arbeit.

\subsection{Grenzen der Aussagekraft}
\label{subsec:diskussion-aussagekraft}

Die messtechnischen Einschränkungen sind in Abschnitt~\ref{sec:grenzen}
aufgeführt. Darüber hinaus ist zu berücksichtigen, wie weit die Ergebnisse
überhaupt verallgemeinert werden können.

Sämtliche Messreihen wurden mit einem Prototyp, an jeweils einem Standort und auf
einem einzigen Kanal erhoben. Insbesondere umfasst der Messaufbau nur eine
empfangende Röhre. Der funktionale Betrieb wurde zwar mit acht gleichzeitig
angesteuerten Röhren erprobt (Abschnitt~\ref{sec:funktionale-tests}) und damit
mit knapp der Hälfte der geforderten Ausbaugröße quantitativ vermessen wurde
dieser Fall jedoch nicht. Dass zusätzliche Röhren die Funkstrecke nicht
zusätzlich belasten, folgt aus der Broadcast-Topologie.

\section{Fazit}
\label{sec:fazit}

Ziel dieser Arbeit war die Entwicklung eines Beleuchtungssystems, das
DMX-Steuersignale drahtlos an mehrere LED-Leuchtröhren überträgt
(Abschnitt~\ref{sec:zielsetzung}). Dieses Ziel wurde erreicht. Das System besteht
aus einer Bridge, die Steuerdaten aus einer bestehenden Lichtinfrastruktur
entgegennimmt oder im Master-Modus selbst erzeugt, sowie aus Leuchtröhren, die
den für ihre Adresse bestimmten Abschnitt des Datenstroms extrahieren und
unmittelbar in Licht umsetzen. Beide Betriebsarten der Bridge und der autarke
Einzelbetrieb der Röhre wurden am aufgebauten System verifiziert.

Die Anforderung an das Echtzeitverhalten wird erfüllt und der Jitter bleibt im gesamten
praxisrelevanten Bereich deutlich unterhalb des Frameabstands. Unter
Idealbedingungen wurde kein einziger Sequenzwert verloren. Die geforderte
Mindestreichweite von 50\,m wird mit freier Sicht deutlich übertroffen und auch
bei vollständig verdecktem Empfänger erreicht. Unter realer, wenn auch moderater
Fremdfunklast bleibt die Übertragung zuverlässig. Die verbleibende Schwachstelle
ist nicht die mittlere Verlustrate, sondern das vereinzelte Auftreten längerer
Unterbrechungen infolge von Kanalkonkurrenz.

Der zweite Zielaspekt die niedrigschwellige Nachbaubarkeit wurde durch die
Beschränkung auf frei verfügbare Komponenten und die vollständige Offenlegung des
Systems umgesetzt. Firmware, Platinenlayouts und Gehäusemodelle liegen unter einer
freien Lizenz vor (\url{https://github.com/nichtrami/led_tube}), sodass sowohl der
Nachbau als auch die Weiterentwicklung durch Dritte möglich ist. Damit adressiert die Arbeit beide Bestandteile der in
Kapitel~\ref{chap:stand-der-technik} identifizierten Lücke: den Integrationsgrad,
indem Funkempfang und Leuchtmittel als Gesamtsystem entworfen wurden, und die
Verfügbarkeit, indem das Ergebnis reproduzierbar dokumentiert und offengelegt ist.

\section{Ausblick}
\label{sec:ausblick}

\subsection{Aus den Messungen abgeleitete Weiterentwicklungen}
\label{subsec:ausblick-messungen}

Die Messreihen weisen mit der Empfangsredundanz auf eine Größe hin, die sich als
Ansatzpunkt für mehrere aufeinander aufbauende Verbesserungen eignet. Sie sinkt
bereits messbar, während der Sequenzverlust noch nahezu bei null liegt, und ist
empfängerseitig ohne Rückkanal bestimmbar
(Abschnitt~\ref{subsec:diskussion-redundanz}).

Die einfachste Nutzung besteht darin, sie sichtbar zu machen. Eine Anzeige der
Verbindungsqualität im Bedienmenü der Röhre erfordert keine Änderung des
Übertragungsverfahrens, da die notwendige Information bereits im Empfangspfad
vorliegt, und gäbe dem Anwender beim Aufbau eine unmittelbare Rückmeldung darüber,
ob eine gewählte Position tragfähig ist.

\subsection{Funktionale Erweiterungen des Gesamtsystems}
\label{subsec:ausblick-funktional}

Neben den aus den Messungen abgeleiteten Verbesserungen bestehen
Erweiterungsmöglichkeiten, die den Funktionsumfang des Systems betreffen.

Der naheliegendste Schritt ist die kabelgebundene Pixelsteuerung über Art-Net. Die
Bridge-Platine sieht hierfür bereits einen über RS-485 geführten Ausgang für die
LED-Datenleitung vor (Abschnitt~\ref{subsec:bridge-pcb}), und die Architektur
überführt sämtliche Eingangsprotokolle ohnehin in ein einheitliches internes
Format (Abschnitt~\ref{subsec:internes-protokoll}). Damit würde das System jene
Betriebsart erschließen, die im Premiumsegment die höhere Auflösung ermöglicht,
und wäre in Aufbauten einsetzbar, in denen eine Datenleitung ohnehin verlegt wird.

Schließlich betrifft eine Weiterentwicklung das Projekt selbst. Da Firmware,
Platinenlayouts und Gehäusemodelle offengelegt sind, entscheidet die Qualität der
Dokumentation darüber, ob ein Nachbau durch Dritte tatsächlich gelingt. Eine
aufbereitete Bau- und Inbetriebnahmeanleitung, die die in
Abschnitt~\ref{subsec:benutzerfreundlichkeit} formulierte Anforderung an den
geringen Einrichtungsaufwand auf den Nachbau überträgt, wäre die konsequente
Fortführung des Open-Source-Ansatzes dieser Arbeit.
